\section{Physics-Structured Inertial Displacement Network}
\label{sec:learning}

\begin{figure}[tb]
\centering
\includegraphics[width=\columnwidth]{architecture.pdf}
\caption{PhysNet. The gyroscope and extrinsic rotation are applied
analytically, and the valid-sample window mean removes the constant
gravity-plus-mean-acceleration component before the network sees the data, so the
trunk only has to supply what is not deterministic. The three heads produce a
velocity stream, a stillness probability that gates that stream,
and a variance; the displacement stream is the integral of the gated velocity
plus the lever-arm term and is anchored at $t_a$ by construction. Supervision is dense: every usable chessboard frame
inside the window contributes, not only the 3-s target.}
\label{fig:architecture}
\end{figure}

\subsection{Physics-Structured Sequence Model}
No external attitude is available at inference here, so our protocol-matched
adaptations of public inertial
networks~\cite{yan2020ronin,liu2020tlio,zeinali2024imunet} consume body-frame
specific force and angular rate and must learn orientation integration,
gravity separation, double integration and the camera--IMU lever arm
implicitly; the original RoNIN and TLIO instead receive gravity-aligned inputs
from an external or filtered attitude. These operations have known analytical structure once the
extrinsic rotation, relative attitude and a gravity estimate are supplied;
they are not inferred from the gyroscope alone. With a few hundred seconds of
reference motion, learning all of them implicitly is wasteful, so PhysNet
(Fig.~\ref{fig:architecture}) computes them explicitly and spends network
capacity only on what is not deterministic: a prior for the initial velocity,
residual bias and motion statistics.

For a pair $(t_a,t_b)$ the input window is $[t_a-c,\,t_b+c]$ resampled to
$f_s=200$~Hz, with $c=2$~s. Gyroscope increments are integrated from the
anchor $t_a$ in both directions after subtracting the mean gyroscope output over
all samples with $\|\tilde{\boldsymbol\omega}_I\|<3^\circ$/s, giving
$\hat{\bf R}_{I_aI_k}$ for every sample $k$. With the convention of
Section~\ref{sec:framework}, the relative camera rotation used below is
${\bf R}_{a\leftarrow k}\equiv{\bf R}_{C_aC_k}
={\bf R}_{CI}\hat{\bf R}_{I_aI_k}{\bf R}_{CI}^{\top}$. The start camera
$C_a$ is also the target frame of
(\ref{eq:target}). Specific force is rotated into $C_a$,
${\bf f}^{C_a}_k={\bf R}_{a\leftarrow k}{\bf R}_{CI}\tilde{\bf f}_{I,k}$.
With ${\cal V}$ the samples inside the recording,
$\hat{\boldsymbol\mu}_{f}=|{\cal V}|^{-1}\sum_{k\in{\cal V}}{\bf f}^{C_a}_k$ is the
window-mean specific force and
$\boldsymbol\ell_k=({\bf f}^{C_a}_k-\hat{\boldsymbol\mu}_{f})\,\mathbb{1}[k\in{\cal V}]$
is the dynamic component used by the network. This is demeaning, not an unbiased gravity compensation:
$\hat{\boldsymbol\mu}_{f}$ contains both the gravity reaction and any nonzero
mean translation acceleration, and both are removed. Its first and second
signed cumulative sums from the anchor are velocity and position pre-integrals
of the demeaned specific force with zero initial velocity. The 24 input
channels are
$\boldsymbol\ell_k$ (3), angular rate mapped from $I_k$ into its instantaneous
camera frame $C_k$ by ${\bf R}_{CI}$ (3), the first two columns of
${\bf R}_{a\leftarrow k}$, the two pre-integrals, the unit vector of
$\hat{\boldsymbol\mu}_{f}$ (opposite to gravity when the mean acceleration is
small), the pair mask, a validity mask, and a normalized time coordinate. All
channels except the two masks and the time coordinate are standardized with
training statistics and clipped to $\pm8$. These channels are computed from
the six-axis IMU stream and timestamps alone; chessboard images and poses are
used only to form training targets and evaluation metrics and are absent at
inference.

\paragraph{Velocity integration and dense supervision.}
Two strided convolutions reduce the sequence to 50~Hz, five dilated residual
convolutions and a bidirectional GRU provide context, and a two-layer head
outputs a velocity stream at 50~Hz that is linearly upsampled to $f_s$,
giving $\hat{\bf v}_k$, the velocity of the IMU sensing point resolved in
$C_a$ (0.42~M parameters). The displacement stream is the discrete integral
\begin{equation}
 \hat{\bf d}_k={\cal I}_{a\rightarrow k}(\hat{\bf v})
 +({\bf R}_{a\leftarrow k}-{\bf I})\,\boldsymbol\lambda ,\qquad
 \hat{\bf d}_a\equiv{\bf 0},
 \label{eq:stream}
\end{equation}
where ${\cal I}_{a\rightarrow k}({\bf x})=f_s^{-1}\sum_{j=a+1}^{k}{\bf x}_j$ for
$k>a$, $-f_s^{-1}\sum_{j=k+1}^{a}{\bf x}_j$ for $k<a$ and ${\bf 0}$ for $k=a$;
$\hat{\bf d}(t)$ denotes linear interpolation of $\hat{\bf d}_k$; and
$\boldsymbol\lambda$ is the lever arm from the IMU origin to the camera
centre expressed in the camera frame
($\boldsymbol\lambda=-{\bf r}_{CI}^{C}$ under the twin convention).
Every usable chessboard frame $m$ inside the
window $[t_a-c,\,t_b+c]$ supervises the stream at its own time,
\begin{equation}
 \begin{aligned}
 {\cal L}={}&\bar\rho\big(\hat{\bf d}(t_b)-\Delta{\bf p}_{ab}\big)
 +\frac{1}{|{\cal M}|}\sum_{m\in{\cal M}}
 \bar\rho\big(\hat{\bf d}(t_m)-\Delta{\bf p}_{am}\big)\\
 &+w_\zeta{\cal L}_{\rm still}+w_\eta{\cal L}_{\rm NLL},
 \end{aligned}
 \label{eq:loss}
\end{equation}
with $\Delta{\bf p}_{am}$ the reference displacement from $t_a$ to $t_m$ by
(\ref{eq:target}), $\bar\rho({\bf r})=\frac13\sum_{q\in\{x,y,z\}}\rho_\beta(r_q)$
the component-averaged smooth-$L_1$ loss with $\beta=0.5$~cm on residuals in
centimetres, and ${\cal M}$ the usable frames of all windows in a minibatch,
over which the dense term is pooled. The last two terms are defined below; the
selected model uses $w_\zeta=1$ and $w_\eta=0$. The first term deliberately
gives $t_b$ one additional unit of weight. A nominal 7-s window contains on average 145 supervised frames instead of one end point, which turns the
regression into a sequence-to-sequence problem and makes overlapping windows
mutually consistent. Augmentation applies per-window bias and white noise to
both sensors, a small random extrinsic rotation, and optionally an exact
time warp by a factor $\gamma$ (demeaned acceleration scaled by $\gamma^2$,
angular rate by $\gamma$) and translation amplitude scaling.

\subsection{Stillness Detection and Velocity Gating}
The unobservable quantity in (\ref{eq:decomp}) is the absolute velocity at the
anchor: an accelerometer only determines velocity increments. A confirmed
stationary phase supplies a zero-velocity boundary condition but remains
subject to detection, attitude and bias errors. A second head on the shared
features outputs a stillness logit $\zeta_k$, supervised by the reference speed
with a dead band---frames slower than 3~mm/s are labelled still ($y_m=1$),
frames faster than 8~mm/s moving ($y_m=0$), and the band between is not
supervised, which keeps boundary noise out of the classifier. The loss
${\cal L}_{\rm still}$ is the binary cross-entropy between $\zeta(t_m)$ and
$y_m$ averaged over the supervised frames. The detected probability then
enters the kinematics before integration,
\begin{equation}
 \hat{\bf v}_k \leftarrow [1-\operatorname{sig}(\zeta_k)]\,\hat{\bf v}_k ,\qquad
 \operatorname{sig}(x)=\frac{1}{1+e^{-x}} .
 \label{eq:gate}
\end{equation}
The reference speed refers to the camera centre whereas the gate acts on the
IMU-point velocity; with a 38-mm lever arm and the median angular rate of
4.3$^\circ$/s the difference is about 3~mm/s, comparable with the still
threshold. This is a soft stillness-conditioned gate rather than a Kalman-style state
reset; it suppresses displacement accumulation when the instrument is likely
to be still. We also
tested a two-pass variant that re-anchors the strapdown velocity at the
detected still samples and feeds the resulting absolute-velocity estimate back
as an input channel; it did not improve beyond the seed spread, alone or
combined with (\ref{eq:gate}), and is not adopted (Table~\ref{tab:ablation}).

A pseudo-measurement also requires an uncertainty statement, not only an
aggregate error, so a third head predicts the per-axis log-variance
$\boldsymbol\eta_k$ of the displacement stream, with
$\eta=\ln(\sigma^2/1~{\rm cm}^2)$ clipped to $[-6,6]$. For displacement residuals
${\bf r}=\hat{\bf d}-\Delta{\bf p}$ in centimetres, the term
${\cal L}_{\rm NLL}=\frac{1}{3|{\cal M}|}\sum_{m}\sum_{q\in\{x,y,z\}}
\frac12[r_{m,q}^2e^{-\eta_{m,q}}+\eta_{m,q}]$ is added to (\ref{eq:loss}). For an
ensemble of $E$ models with end points $\hat d_q^{(e)}$, mean $\bar d_q$ and
variances $\sigma_q^{2(e)}=e^{\eta_q^{(e)}}$, the reported standard uncertainty
per axis is
$u_q=\kappa\,[\bar\sigma_q^2+s_{{\rm ens},q}^2]^{1/2}$ with
$\bar\sigma_q^2=E^{-1}\sum_e\sigma_q^{2(e)}$ and
$s_{{\rm ens},q}^2=E^{-1}\sum_e(\hat d_q^{(e)}-\bar d_q)^2$. The single scale
$\kappa$ is the RMS standardized residual on validation~\cite{laves2020well},
and the result is assessed by interval
coverage~\cite{kendall2017uncertainties,lakshminarayanan2017simple}.

\subsection{Protocols and Trajectory Formation}
Three protocols are used. The \emph{fixed split} fits on six training
recordings and uses the two validation recordings for early stopping and all
design choices. \emph{Leave-one-session-out} holds out each of the eight
non-test recordings in turn while fitting on the other seven plus the
rotation-rich rigidity recordings (\emph{extra-train}), which scores every
non-test second exactly once and yields a session-level error distribution;
the fold models are then averaged and applied to the held-out test
recording after fold training. \emph{Synthetic direct} trains on the twin only and is evaluated
without real updates. Two blends of PhysNet and the strongest public
architecture are reported: the \emph{equal blend} averages the two ensemble
predictions, and the \emph{calibrated blend} uses a convex weight $w$ and a
global scale $\alpha$, $\alpha[w\hat{\bf d}_{\rm P}+(1-w)\hat{\bf d}_{\rm I}]$,
where $\hat{\bf d}_{\rm P}$ and $\hat{\bf d}_{\rm I}$ are the PhysNet and IMUNet
ensemble end points; $w\in\{0,0.05,\dots,1\}$ and $\alpha\in[0.6,1.2]$ (step
0.025) minimize the mean error of the out-of-fold cross-validation
predictions.

For trajectory evaluation, displacement pseudo-measurements are converted to
board-frame edges using the reference orientation when translation is assessed
in isolation, and a sparse least-squares pose graph solves
\begin{equation}
 \begin{aligned}
 &\min_{\{\hat{\bf p}_i\}}\sum_{(i,j)\in{\cal E}}
 \alpha_{ij}^{2}\|(\hat{\bf p}_j-\hat{\bf p}_i)-{\bf R}_{BC}(t_i)
 \widehat{\Delta{\bf p}}_{ij}\|^2\\
 &\text{s.t.}\quad\hat{\bf p}_{i_0}={\bf p}_{BC}(t_{i_0}),
 \end{aligned}
\end{equation}
where $\hat{\bf p}_i$ is the position of node $i$ in $B$, ${\cal E}$ is the edge
set, and $i_0$ is the first node of each connected component, which removes the
translation gauge per component. End-point regressors provide one edge per
window with $\alpha_{ij}=1$; for the dense variant the PhysNet stream
(\ref{eq:stream}) replaces it by an edge from the anchor to every usable frame
in the window with $\alpha_{ij}=e^{-|t_j-t_i|/(1.5\,{\rm s})}$. This experiment measures temporal
error correlation; it is not reported as pure-IMU pose because reference
orientation is supplied.

A separate relative-pose diagnostic removes the orientation assistance. The
first reference attitude fixes the rotational gauge, gyroscope increments
propagate orientation thereafter, and the predicted local translations are
rotated by that inertial attitude before the same pose-graph solve, including
its per-component position anchors. Orientation is gyro-integrated rather than
learned, and the result is not an absolute-pose estimate.
